\section{De la puerta de entrada publicitaria al trabajo intelectual: una revaloración del modelo de negocio del Agent}

La discusión sobre el modelo de negocio empieza con una comparación muy clásica: los grandes modelos de negocio de la era de internet son sobre todo publicidad, suscripción y comisión por transacción. La publicidad vende atención, la suscripción vende eficiencia de herramienta, la comisión vende intermediación de transacciones. Cuando el Agent funciona como puerta de entrada, desafía directamente a la publicidad, porque el sistema publicitario quiere que el usuario permanezca, navegue y haga clic, mientras que un buen Agent debería precisamente reducir el tiempo que el usuario permanece, y hacer por él el filtrado, la comparación y la ejecución.

El juicio de Shang Yanyi (尚晏仪) resulta muy revelador: lo que completa internet es la conexión de información, y el gasto publicitario corresponde a la distribución de información; lo que produce el Agent es trabajo intelectual, y en la sociedad humana el trabajo intelectual suele pagarse mediante salario, cuota de servicio, comisión de intermediación u honorario. Por eso ella confía más en la suscripción y la comisión que en forzar al Agent de vuelta a la lógica publicitaria. Si el Agent en verdad funciona como secretario, asesor o representante, su valor comercial debería provenir de ahorrar tiempo, reducir el costo de decisión, intermediar transacciones y asumir el resultado de la acción, y no de hacer que el usuario vea el contenido unos segundos más.

El escenario de los viajes funciona en su práctica precisamente porque en él existe de manera natural el papel histórico del travel agent. El usuario entiende qué significa un «representante» y también acepta pagar por planeación, filtrado, transacción y servicio. Los viajes además tienen un ciclo de transacción claro: hotel, vuelo, ruta, visa y servicios de itinerario pueden conectarse a una comisión o a una cuota de servicio. Esto es distinto de una «superpuerta de entrada» generalizada, que parece tener más espacio pero cuyo reward, cumplimiento y ciclo comercial son más difusos.

\begin{importantbox}{El Agent se parece más a una fuerza de trabajo que a un espacio publicitario}
Si el valor de un sistema consiste en pensar, filtrar, actuar y asumir el costo del proceso en lugar del usuario, entonces se acerca más al trabajo intelectual. La lógica publicitaria exige ocupar la atención del usuario; la lógica del Agent exige liberar la atención del usuario. Las dos no son incompatibles, pero el centro de negocio por defecto es distinto.
\end{importantbox}

La discusión sobre el teléfono Doubao lleva este asunto a un nivel más alto. La ventaja del Agent a nivel de teléfono es muy clara: el teléfono conoce la dirección, las preferencias, las cuentas, las conversaciones, la agenda y los hábitos de vida del usuario, por lo que tiene la mejor oportunidad de convertirse en secretario de vida. Pero al mismo tiempo toca la estructura de intereses existente en la distribución de aplicaciones, la recomendación publicitaria y las puertas de entrada a las transacciones. Dicho de otro modo, el Agent de teléfono es técnicamente muy razonable, pero comercialmente cambiará el sustento de muchas personas.

Aquí hay un detalle que vale la pena retener. Shang Yanyi dice que si el usuario le dice directamente al Agent «quiero ir a una isla», esto es más contundente que un sistema de publicidad de video corto que observa que el usuario se detiene unos segundos más en videos de islas. Porque lo primero es una intención explícita, mientras que lo segundo es solo una inferencia de conducta. Pero la intención explícita no equivale automáticamente a éxito comercial. El Agent todavía tiene que estructurar la intención, encontrar la oferta, comparar opciones, construir confianza y completar la transacción. El sistema publicitario es hábil para convertir conductas ambiguas en recomendaciones; el sistema de Agent, en cambio, tiene que convertir lenguaje explícito en acción confiable.

\begin{knowledgebox}{De la economía de la atención a la economía del representante}
\begin{tabularx}{\textwidth}{>{\raggedright\arraybackslash}p{0.2\textwidth} Y Y}
\toprule
Modelo & Origen del valor & Relación con el Agent \\
\midrule
Publicidad & Permanencia, navegación, clic y exposición del usuario. & Un buen Agent reduce la permanencia, por lo que hay un conflicto natural. \\
Suscripción & El usuario paga por eficiencia, capacidad de herramienta y confiabilidad. & Apto para escenarios de alta frecuencia, explícitos y de mejora de eficiencia sostenible. \\
Comisión & La plataforma intermedia la transacción y participa en el cumplimiento. & Apto para viajes, compras, comercio electrónico, reclutamiento y otros escenarios con ciclo de transacción. \\
Servicio de representación & El usuario paga por trabajo intelectual, filtrado, ejecución y límites de responsabilidad. & La forma de largo plazo más cercana a un Agent fuerte. \\
\bottomrule
\end{tabularx}
\end{knowledgebox}

\subsection{Resumen de la sección}

El producto Agent no puede simplemente replicar el modelo publicitario de internet. Su valor se parece más al trabajo intelectual, al servicio de representación y a la intermediación de transacciones. Los viajes son adecuados como caso precisamente porque tienen tanto un papel tradicional de representante como un ciclo de transacción claro. El potencial del Agent de teléfono proviene del contexto personal, pero lo verdaderamente difícil es convertir ese contexto en acción confiable, mientras se manejan los conflictos comerciales con el ecosistema existente.
